\documentclass[11pt]{article}
\usepackage[utf8]{inputenc}
\usepackage[T1]{fontenc}
\usepackage{amsmath,amssymb}
\usepackage{booktabs}
\usepackage{tabularx}
\usepackage[authoryear,round]{natbib}
\usepackage[margin=1in]{geometry}
\usepackage{microtype}
\setlength{\emergencystretch}{3em}
\usepackage[hidelinks]{hyperref}

% Calibration numbers: see the Zenodo deposit.

\newcolumntype{L}{>{\raggedright\arraybackslash}X}
\newcommand{\chk}[1]{\texttt{#1}}
\newcommand{\term}[1]{\item[\textbf{#1}]}
\newcommand{\lvl}[2]{\subsubsection*{ARL #1: #2}}

\title{The Alpha Readiness Level: A Calibrated Evidence Scale\\ for Trading-Strategy Backtests}
\author{Rostyslav Shcherbyna\\ Independent researcher, Zaporizhzhia, Ukraine\\
ORCID: \href{https://orcid.org/0009-0000-3773-5789}{0009-0000-3773-5789}}
\date{28 November 2026}

\begin{document}
\maketitle

\begin{abstract}
A backtest shows how a trading strategy would have performed on past data. Its headline number, usually the Sharpe
ratio, says how good the result looked, but not how much evidence stands behind it. This paper proposes the Alpha
Readiness Level (ARL), an ordinal scale from 1 to 9 modelled on technology readiness levels. Each level is a named
claim, such as ``it earns after costs'' or ``it survives the search'', with requirements that software can check. A
level counts only if every level below it is met. I calibrate the scale on cases whose truth is known: deliberate
traps, strategies with no edge, and strategies with an edge of controlled size. Development on 350,000 cases exposed
two structural failures: the best of 100 or 1,000 random strategies, reported as a single attempt, reached level 4 in
93\,\% of cases, and no strategy with a real single-asset edge could reach level 5. I then froze the revised rules,
pre-registered four predictions and ran 350,000 new cases; all four were met. Strategies with no edge exceeded their
maximum level in 0.82\,\% of cases (95\,\% interval 0.76--0.89; 1.81\,\% on real market bars), and real single-asset
and cross-sectional edges reached level 6 in 29.0\,\% and 52.8\,\% of cases. The share of traps promoted too far
depends on what the user declares: 0.80\,\% when traps that hide trials keep no record of them and cross-sectional traps
are given no added universe file, but 8.29\,\% (8.10--8.48) when every trap that can take them is given a new trial
record and, if cross-sectional, a universe declaration. A new record that omits the search lets the hidden-trials trap
reach level 4 or above in 46.9\,\% of cases. An ARL level is therefore evidence on the condition that the declarations
behind it are truthful. The rules and the verification software are open source.
\end{abstract}

\noindent\textbf{Keywords:} backtesting; overfitting; multiple testing; deflated Sharpe ratio; pre-registration;
research protocol; technology readiness level; quantitative finance

\subsection*{Summary for the general reader}

\paragraph{The problem.} A computer can test thousands of trading rules on past prices in minutes. Some of them will
look excellent by pure luck. A reader who sees only the best result cannot tell luck from skill.

\paragraph{The proposal.} ARL answers a narrower and more honest question than ``is this strategy good?''. It asks:
\emph{how far has the evidence been taken?} Level 1 means that an idea was tried on usable data. Level 6 means that the
result still holds after accounting for the number of attempts. Level 9 means that it was confirmed with real orders
and real money.

\paragraph{The main findings.}
\begin{enumerate}
\item Hidden attempts cannot be detected from one result. They have to be recorded while the research is done.
\item With the final rules, strategies with no edge were rarely promoted too far: about 1 case in 120.
\item Real but weak edges need long samples. A low level often means ``not enough data yet'', not ``no edge''.
\item The scale believes what the user declares. A false or empty declaration can raise the level, and in such cases the
rate of wrongly promoted traps rose from 0.80\,\% to 8.29\,\%.
\item A level should therefore always be read together with the evidence that supports it.
\end{enumerate}

\paragraph{What ARL is not.} It is not a forecast of profit, not a rating of a strategy's quality, and not investment
advice. It is also not yet an accepted industry standard: it is an open proposal that others can apply, test and
criticise.

\section{Introduction}

Most strategies that look profitable in a backtest do not survive contact with new data. The reasons are well
documented: errors that let the code see the future, costs that were left out, and above all the selection of the best
of many variants \citep{bailey2014deflated,harvey2016cross,arnott2019backtesting}. Each failure has a remedy in the
literature, but the remedies are scattered, and a reader of a backtest report rarely knows which of them were applied.

Consider two reports that both state an annual Sharpe ratio of 1.5. The first strategy was fixed in advance, tested
once, charged realistic costs and later traded with real money. The second is the best of 1,000 variants, tested
without costs. The headline number is the same; the evidence is not. A single Sharpe ratio hides this difference.

Engineering faced a similar problem with new technologies and answered it with technology readiness levels
\citep{mankins1995trl}: an ordered scale in which each level names a concrete demonstration, from a principle observed
to a system proven in operation. The scale does not say that a technology is good; it says how far the evidence goes.
This paper proposes the same idea for trading strategies. The Alpha Readiness Level (ARL) orders the questions that a
sceptical reviewer asks of a backtest, attaches requirements that software can check, and reports the highest level
whose requirements are met together with those of every level below it.

A scale of this kind is useful only if it does not promote what it should not. I therefore treat its thresholds as
hypotheses and test them on cases whose truth is known, most of them synthetic: traps with a known maximum level,
strategies without any edge, and edges of controlled size planted in the data. This calibration did more than tune
thresholds. It exposed two failures of structure that no threshold could repair, and their remedies are the main
practical result of the paper.

The contributions are:
\begin{enumerate}
\item the scale itself, with nine levels, ceilings, blockers, a depth sub-scale and rules of honest use
(Section~\ref{sec:scale}), illustrated by worked examples (Section~\ref{sec:example});
\item a calibration design with three populations of cases and two verification tracks (Section~\ref{sec:calib});
\item development results for three versions of the rules, including two structural repairs and their side effects
(Section~\ref{sec:results});
\item a pre-registered confirmatory run of the frozen final rules on new cases (Section~\ref{sec:confirm});
\item a recommended way of reporting an ARL level (Section~\ref{sec:use}) and an open, reproducible implementation in
which the rules are data rather than program code (Section~\ref{sec:impl}).
\end{enumerate}

\section{Related work}

\paragraph{Selection and multiple testing.} The deflated Sharpe ratio corrects a Sharpe ratio for the number of trials
and for returns that are not normally distributed \citep{bailey2014deflated}. The probability of backtest overfitting
estimates how often the configuration that is best in sample underperforms out of sample \citep{bailey2017pbo}.
\citet{harvey2016cross} raise the t-statistic hurdle for new factors to about 3. Earlier, \citet{white2000reality} and
\citet{hansen2005spa} gave tests of the best of many models against a benchmark. All of these need the number of trials,
which the analyst reports.

\paragraph{Statistics of the Sharpe ratio and of track records.} \citet{lo2002sharpe} derives the sampling distribution
of the Sharpe ratio under serial correlation; \citet{getmansky2004} show how smoothed returns inflate it.
\citet{bailey2014pseudo} derive the minimum backtest length for a given number of trials.

\paragraph{Research protocols.} \citet{arnott2019backtesting} list rules for a credible backtest. Pre-registration,
standard in clinical research and promoted for science at large \citep{nosek2018prereg}, fixes the hypothesis before the
data are seen. \citet{mclean2016} find that published anomalies lose much of their return after publication.

\paragraph{Readiness and disclosure.} Technology readiness levels have been adapted to machine-learning systems
\citep{lavin2022mltrl}, and \citet{fabozzi2018honest} propose a template for disclosing the number of tests behind a
backtest.

ARL does not replace these tools. It arranges them into a sequence of claims and states which evidence each claim
needs. To my knowledge, no graded evidence scale for trading-strategy backtests has been published together with a
calibration on cases of known truth.

\section{Concepts and definitions}\label{sec:defs}

This section defines the terms used in the rest of the paper. Readers familiar with quantitative research may skip to
Section~\ref{sec:scale}.

\subsection{From backtest to live trading}\label{sec:modes}

A strategy can be exercised in several environments. They differ in who decides whether an order is filled, at what
price, and whether money is at stake. A \emph{fill} is the execution of an order: its quantity, price and time.

\begin{description}
\term{Backtest.} The strategy is replayed on historical data. No order is sent anywhere. Fills are computed from
assumptions, for example ``buy at the next open and pay 10 basis points''.
\term{Simulated fills.} Any fills produced by a model of the market, including a backtest or a replay. They can only be
as realistic as the model.
\term{Paper trading.} The strategy runs on current market data as time passes, but its orders are only written down.
Software decides how they would have been filled. No money is involved. Paper trading shows that the strategy can
produce decisions in time on data it has never seen; it does not show that the market would have filled them.
\term{Broker demo account.} A practice account that a broker offers with virtual money. Orders go to the broker's demo
system, which decides the fills by its own rules. It exercises the connection to the broker and the order handling.
Prices may resemble live prices, but no real counterparty takes the other side, so slippage, partial fills and
rejections can differ from live trading.
\term{Exchange testnet.} A separate test copy of an exchange, with test balances that have no value. An example is the
Binance Spot Test Network\footnote{\url{https://developers.binance.com/en/docs/products/spot/testnet/general-info},
accessed 28 November 2026.}: its funds are virtual and cannot be transferred out, its trading interface mirrors that of
the real exchange, and it is reset to a blank state from time to time. Orders meet only the orders and liquidity of the
test system. It proves that the code can place, change and cancel orders; it does not prove that those orders could be
executed at the same prices and sizes on the real market.
\term{Live trading.} Orders with real money reach the real market. Fills, fees, queue position, market impact,
rejections and outages are real, and so are the losses.
\end{description}

\begin{table}[ht]
\centering\small
\begin{tabularx}{\linewidth}{@{}p{0.17\linewidth}LLp{0.14\linewidth}p{0.12\linewidth}@{}}
\toprule
Environment & Market data & Who decides the fill & Money & Highest ARL \\
\midrule
Backtest & historical & the researcher's model & none & 8 \\
Paper trading & current & software of the user & none & 8 \\
Broker demo & current, from the broker & the broker's demo system & virtual & 8 \\
Exchange testnet & the test system's own & the test exchange & test assets & 8 \\
Live trading & current, real & the real market & real & 9 \\
\bottomrule
\end{tabularx}
\caption{Environments in which a strategy can be exercised. The verifier accepts four declared sources of forward
results: \chk{real\_fills}, \chk{paper}, \chk{testnet} and \chk{simulated}; a broker demo account is not real fills and
is declared as paper or simulated. Only real fills can support level 9.}\label{tab:modes}
\end{table}

\subsection{Ways in which a backtest misleads}

\begin{description}
\term{Look-ahead bias.} The strategy uses information that was not available at the time of the decision. Example:
deciding today's position from tomorrow's closing price, or standardising a series with the mean of the whole sample.
\term{Survivorship bias.} The assets tested are those that still exist. Example: testing today's index members over six
years ignores the companies that were removed, often after large losses.
\term{Point-in-time universe.} The remedy: a \emph{membership file} that states for every asset the dates on which it
belonged to the tradable set.
\term{Costs, spread and slippage.} A trade pays commission and the difference between buying and selling prices, and is
often filled worse than expected. Costs are given in \emph{basis points} (bps); 1~bps is 0.01\,\%.
\term{Selection bias (multiple testing).} The more variants are tried, the better the best one looks by chance. Each
variant tried is a \emph{trial}. A person who tries 1,000 random rules and reports only the best one reports luck.
\term{Overfitting.} Parameters are tuned to the accidents of one sample and fail on new data.
\term{In sample and out of sample.} Data used to choose a strategy are in sample; data kept away from every choice are
out of sample. A \emph{sealed holdout} is an out-of-sample period that is frozen in advance and opened once.
\end{description}

\subsection{Performance and statistical evidence}

\begin{description}
\term{Sharpe ratio.} Mean return divided by its standard deviation, usually per year:
$\mathrm{SR}=\sqrt{P}\,\bar r/s$, where $P$ is the number of periods in a year. Example: daily returns with mean
0.05\,\% and standard deviation 1\,\%, on 365 days a year, give $\mathrm{SR}\approx 0.96$. A \emph{net} Sharpe ratio is
computed after costs.
\term{Deflated Sharpe ratio (DSR).} A probability between 0 and 1 that the true Sharpe ratio exceeds the best value
expected from the stated number of skill-less trials, allowing for sample length and for skewed, fat-tailed returns
\citep{bailey2014deflated}. The more trials are declared, the harder it is to pass.
\term{Harvey--Liu--Zhu (HLZ) hurdle.} The requirement that the t-statistic of the mean return be at least 3 rather
than the conventional 2 \citep{harvey2016cross}.
\term{Minimum backtest length.} The history needed so that the best of $N$ random strategies would no longer be expected
to reach the observed annual Sharpe ratio, approximately $2\ln N/\mathrm{SR}^2$ years \citep{bailey2014pseudo}. Example:
for $N=100$ and $\mathrm{SR}=1$ this is about 9.2 years.
\term{Probability of backtest overfitting (PBO).} The share of ways of splitting the data in which the configuration
that is best in one half ranks below the median in the other \citep{bailey2017pbo}.
\term{p-value.} The chance of seeing a result at least as good if there were no effect. Small values speak against luck.
\term{Information coefficient (IC).} For strategies that rank many assets: on each date, the rank correlation between
the scores given to the assets and their subsequent returns, averaged over dates. Example: a mean IC of 0.03 says that
higher-scored assets tend, weakly, to do better. The verifier also reports its t-statistic.
\term{Placebo score and placebo test.} A \emph{placebo score} is a fake signal made by shuffling the real scores among
the assets on each date, so that any link to returns is destroyed while everything else is kept. The real strategy is
compared with many placebos (200 by default, 100 in the calibration). The p-value is the share of placebos that do at
least as well. If shuffled scores earn as much as real ones, the profit comes from the way the portfolio is built, not
from the signal.
\term{Timing test.} The corresponding test for a single asset: the series of positions is shifted in time by many random
offsets (200 by default). A strategy that times the market should beat nearly all shifted copies of itself.
\term{Alpha, beta, factor exposure.} \emph{Beta} is the part of the return explained by holding the market or a known
\emph{factor} such as momentum. \emph{Residual alpha} is what remains after that part is removed.
\term{Hidden short volatility.} A payoff that earns small steady gains and occasionally loses heavily when markets move
sharply, like selling insurance.
\term{Break-even cost.} The cost per side at which the strategy's total profit becomes zero. Example: if the assumed cost
is 10~bps and the break-even cost is 45~bps, the margin is 4.5 times.
\term{Market impact and capacity.} Large orders move prices against the trader. \emph{Capacity} is how much money the
strategy can run before this, or the share of traded volume it needs, becomes too large.
\term{Book and marginal Sharpe ratio.} The \emph{book} is the set of strategies already run. The marginal Sharpe ratio is
the improvement in the book's Sharpe ratio when the new strategy is added.
\end{description}

\subsection{Terms of verification and calibration}

\begin{description}
\term{Check and status.} A check is one automated test, named in this paper in \chk{typewriter} type. It returns
\emph{pass}, \emph{info} (informative, nothing wrong found), \emph{warn}, \emph{fail} or \emph{skip} (not run).
\term{Certificate.} The record of one verification: checks, numbers, settings, rules version and a hash (fingerprint)
of the data.
\term{Trial ledger.} A file in which every variant tried is appended as it is run. Each entry is chained to the previous
one by a hash, so later edits are visible.
\term{Pre-registration.} A time-stamped statement of the hypothesis and method, made before the results are seen.
\term{Trap.} A test case built with a known flaw and therefore a known maximum level.
\term{False promotion.} A trap or a no-edge strategy that receives a level above its known maximum.
\term{Power.} The share of strategies with a real edge that reach a given level.
\term{Wilson 95\,\% interval.} A range that expresses the sampling uncertainty of a share estimated from a finite number
of cases \citep{wilson1927}.
\end{description}

\section{The scale}\label{sec:scale}

\subsection{How a level is computed}

\begin{enumerate}
\item The verifier runs the checks for which the user supplied the needed inputs.
\item Each level lists the checks it requires. Some levels accept alternatives, because different evidence answers the
same question for different kinds of strategy.
\item A level is \emph{met} if its required checks ran and returned an accepted status.
\item The ARL is the highest level $L$ such that level $L$ and every level below it are met.
\item Ceilings and blockers (Section~\ref{sec:caps}) are then applied.
\item A required check that did not run leaves its level \emph{not assessed}. Missing evidence is never counted as
passed, and the count stops there.
\end{enumerate}

A strategy with a highly significant signal but a backtest that looks into the future is therefore at level 1: later
evidence cannot compensate for a broken foundation.

\begin{table}[ht]
\centering\small
\begin{tabularx}{\linewidth}{@{}rp{0.27\linewidth}L@{}}
\toprule
L & Claim & Question answered \\
\midrule
1 & Idea on clean data & Was the idea tried on usable data? \\
2 & The backtest is not broken & Does the backtest use only information available at the time? \\
3 & It earns after costs & Is there still a profit after realistic costs? \\
4 & The signal predicts & Does the signal beat chance, given the attempts made? \\
5 & It is new alpha & Is it more than market or factor exposure in disguise? \\
6 & It survives the search & Does it stand up to a strict correction for the search? \\
7 & Tradable and robust & Could it be traded, and does it hold under changed conditions? \\
8 & It adds to your book & Does it improve what is already run? \\
9 & Confirmed live & Did real orders confirm it? \\
\bottomrule
\end{tabularx}
\caption{The nine levels of \texttt{arl-rules/2}.}\label{tab:levels}
\end{table}

\subsection{The nine levels in detail}\label{sec:levels}

The thresholds below are those of \texttt{arl-rules/2} in the open-source verifier Monte-Neo, release
\texttt{v0.57.0}. The examples are illustrations; Section~\ref{sec:example} gives executed ones.

\lvl{1}{Idea on clean data}
Required, one of:
\begin{itemize}
\item \chk{data\_integrity} is not \emph{fail}: the series is not broken, for example dates do not run backwards;
\item an information coefficient (\chk{ic}) could be computed.
\end{itemize}
\emph{Passes:} a daily price table with increasing dates and valid prices. \emph{Stops:} a table whose dates jump
backwards. Level 1 says nothing about profit.

\lvl{2}{The backtest is not broken}
Required, one of:
\begin{itemize}
\item \chk{lookahead\_truncation} passes: the signal is recomputed on data cut off at several points, and past positions
do not change when the future is removed;
\item \chk{universe\_lookahead} is not \emph{fail}: no position was held in an asset outside its membership dates.
\end{itemize}
Any failed look-ahead check caps the level at 1. \emph{Passes:} a moving average of past closes. \emph{Stops:} a
five-day average centred on today, which includes the next two days.

\lvl{3}{It earns after costs}
Required, one of:
\begin{itemize}
\item single-asset strategies: \chk{net\_profitability} passes (total return after costs above zero) and
\chk{cost\_margin} passes (break-even cost at least 3 times the modelled cost per side);
\item portfolios of many assets: \chk{market\_impact} is not \emph{fail} (positive annual return after commission and
price impact at the stated capital) and \chk{cost\_sensitivity} is not \emph{fail} (positive at the modelled costs).
\end{itemize}
\emph{Passes:} costs of 10~bps and break-even at 45~bps. \emph{Stops:} costs of 10~bps and break-even at 22~bps; the
strategy is profitable, but its margin is too thin.

\lvl{4}{The signal predicts}
Required, all of:
\begin{itemize}
\item evidence of prediction, either (a) a positive mean information coefficient with $|t|\ge 2$ (\chk{ic}) and a
placebo test with $p\le 0.10$ (\chk{signal\_placebo}), or (b) a timing test with $p\le 0.05$
(\chk{timing\_significance});
\item a deflated Sharpe ratio of at least 0.8, computed with at least the project's number of trials;
\item \emph{trial evidence}: an intact trial ledger, trials imported from an experiment tracker, a pre-registration made
before the first verification of the same code, or a research-process record. Without it the level is capped at 3.
\end{itemize}
\emph{Passes:} a signal that beats its placebos, with its attempts on record. \emph{Stops:} a good-looking result with
no record of how many variants preceded it.

\lvl{5}{It is new alpha}
Required, both of:
\begin{itemize}
\item \chk{residual\_alpha} passes: after removing exposure to the market or known factors, the remaining return has a
t-statistic of at least 3;
\item \chk{hidden\_short\_volatility} passes: no signs of a payoff that sells insurance.
\end{itemize}
No blocker of Section~\ref{sec:caps} may have warned or failed. \emph{Passes:} a return that remains after the market
and momentum are removed. \emph{Stops:} a ``new'' signal that is ordinary 12-month momentum under another name.

\lvl{6}{It survives the search}
Required, one of:
\begin{itemize}
\item a deflated Sharpe ratio of at least 0.95 (\chk{deflated\_sharpe});
\item the HLZ hurdle $t\ge 3$ (\chk{hlz\_hurdle}) together with a history not shorter than half the minimum backtest
length when ten or more trials are declared (\chk{min\_backtest\_length} not \emph{fail}).
\end{itemize}
\emph{Passes:} ten years of a strong, stable edge with 20 declared trials. \emph{Stops:} one year of the same edge,
which is too short to exclude luck.

\lvl{7}{Tradable and robust}
Required, all three groups:
\begin{itemize}
\item \emph{tradability:} \chk{market\_impact} or \chk{capacity} was assessed and did not fail;
\item \emph{robustness}, either (a) the strategy beats resampled artificial markets
(\chk{synthetic\_market\_null}, $p\le 0.10$), holds on random subsets of the assets (\chk{universe\_subsampling}, not
\emph{fail}) and does not depend on the start date (\chk{start\_date\_sensitivity}), or (b) its profit does not come
from a single period (\chk{period\_consistency}) and survives a delay of one bar (\chk{delay\_sensitivity});
\item \emph{execution margin:} no strong dependence on the day of rebalancing (\chk{rebalance\_timing\_luck}) or an
adequate cost margin (\chk{cost\_margin}).
\end{itemize}
\emph{Passes:} profitable in several periods, still profitable when orders arrive one bar late, with room for costs.
\emph{Stops:} a profit that disappears with a one-bar delay.

\lvl{8}{It adds to your book}
Required: \chk{marginal\_sharpe} passes: adding the strategy raises the book's Sharpe ratio by more than 0.05, and the
90\,\% bootstrap interval of the gain excludes zero. \emph{Passes:} a strategy that earns when the existing book does
not. \emph{Stops:} a near copy of a strategy already held.

\lvl{9}{Confirmed live}
Required, all of:
\begin{itemize}
\item \chk{live\_vs\_backtest} passes: over at least 20 common periods, live and backtest returns have a correlation of
at least 0.5 and live results do not trail the backtest significantly;
\item \chk{alpha\_death\_monitor} passes: sequential tests do not signal that the edge has vanished;
\item \chk{live\_source} passes: the data are real fills, with at least 60 live bars (more for weaker edges) and at
least 30 fills whose slippage is within 1.5 times the modelled value.
\end{itemize}
\emph{Passes:} months of small real orders that behave like the backtest. \emph{Stops:} an excellent testnet or demo
record, which is capped at 8.

\subsection{Ceilings and blockers}\label{sec:caps}

Some findings limit the level whatever else has passed.

\paragraph{Ceilings} set a maximum level:
\begin{itemize}
\item a failed look-ahead check: at most 1;
\item no trial evidence: at most 3;
\item a universe that may consist only of today's survivors (check \chk{survivorship}): at most 4;
\item paper, testnet or simulated fills instead of real ones: at most 8.
\end{itemize}

\paragraph{Blockers} stop the strategy below level 5 when the check ran and warned or failed. A blocker that did not run
does not block.
\begin{itemize}
\item \chk{alpha\_after\_beta}: the return is not shown to be more than exposure to a benchmark;
\item \chk{universe\_subsampling}: the result rests on a few assets;
\item \chk{noise\_injection}: the result does not survive small noise in the inputs;
\item \chk{short\_feasibility}: the short side could not realistically be borrowed or traded;
\item \chk{residual\_vs\_book}: the return is explained by the existing book;
\item \chk{cross\_market\_close\_alignment}: the signal uses closing prices of markets that close after the trade;
\item \chk{post\_discovery\_decay}: the edge disappears after the date of its discovery;
\item \chk{claim\_consistency}: the claimed numbers are better than the verified ones;
\item \chk{universe\_lookahead}: the universe looks like today's members and no membership file was given. For a fixed
universe, such as a set of currency pairs, the file lists every asset for the whole period.
\end{itemize}

\subsection{Depth: how strong is the evidence at that level?}\label{sec:depth}

The level says which claim holds. \emph{Depth} says how firmly it holds, on a count from 0 to 4. At the level reached,
one point is given for each criterion that is met:
\begin{enumerate}
\item \textbf{D1, margin:} the main result is clearly beyond the threshold;
\item \textbf{D2, breadth:} it repeats across periods, assets, regimes or probes;
\item \textbf{D3, sample:} there are enough years, observations or trades;
\item \textbf{D4, out of sample:} it holds again on data that played no part in any choice.
\end{enumerate}

A result is written as ``ARL 6, depth 3/4''. Four rules govern depth:
\begin{itemize}
\item depth never raises the level: ARL 6 with depth 4 remains below ARL 7 with depth 0;
\item a criterion whose evidence was not produced scores 0 and is shown as \emph{not assessed}, with the command that
would produce it;
\item the sealed holdout counts only if it was frozen and opened once; a second opening sets D4 to zero for every level,
permanently;
\item at levels 4 to 6 the statistics are adjusted for the number of trials of the whole project.
\end{itemize}

Table~\ref{tab:depth} lists the criteria of every level.

\begin{table}[p]
\centering\footnotesize
\begin{tabularx}{\linewidth}{@{}rLLLL@{}}
\toprule
L & D1 margin & D2 breadth & D3 sample & D4 out of sample \\
\midrule
1 & no data warning at all & two independent data sources agree & at least 3 years of bars & a fresh download gives identical hashes \\
\addlinespace
2 & all four look-ahead probes pass & truncation and perturbation probes pass at 12 or more cut points & truncation probe passes at 20 or more cut points & probes pass again on the sealed holdout \\
\addlinespace
3 & break-even cost at least 5 times the modelled cost & profitable in at least 5 of 6 equal periods & at least 100 closed trades & net profit on the sealed holdout \\
\addlinespace
4 & placebo or timing $p$-value, multiplied by the number of trials, below 0.01 & positive in at least 70\,\% of periods or asset subsets & at least 1,000 observations & significant on the sealed holdout \\
\addlinespace
5 & residual alpha $t\ge 3$ after adjustment for the trials & alpha remains after all factors, including the market & at least 5 years & residual alpha positive on the sealed holdout \\
\addlinespace
6 & deflated Sharpe ratio at least 0.99 & probability of backtest overfitting below 0.2 & trial count checked against the ledger & deflated Sharpe at least 0.95 on the sealed holdout \\
\addlinespace
7 & profit at doubled costs and with a 1-bar delay & still passes across regimes, rebalance days and (if run) subsets and start dates & capacity at least 10 times the planned capital & survives the holdout under realistic fills \\
\addlinespace
8 & interval of the marginal Sharpe gain excludes zero & positive in sub-periods and stress periods & at least 3 years of overlap with the book & positive on forward data \\
\addlinespace
9 & live Sharpe within the backtest's confidence interval & seen in at least 2 market conditions & at least 6 months and 100 live fills & real fills with a custody statement \\
\bottomrule
\end{tabularx}
\caption{Depth criteria of \texttt{arl-rules/2}. A criterion is scored only when the corresponding evidence exists. The
look-ahead probes of level 2 are truncation, perturbation of future data, a static scan of the source code and a test
that past positions are not redrawn. A \emph{custody statement} is a record from the broker or custodian that supports
the live results.}\label{tab:depth}
\end{table}

\paragraph{Reading depth.} Three examples, the first two taken from the executed cases of
Section~\ref{sec:example}:
\begin{itemize}
\item \emph{ARL 3, depth 2/4} (case B). The strategy earns after costs. D2 and D3 are met: it was profitable in 5 of 6
periods and had more than 100 trades. D1 is not met: its break-even cost is about 3 times the modelled cost, enough for
level 3 but short of 5 times. D4 was not assessed, because no holdout was given.
\item \emph{ARL 7, depth 1/4} (case E). The strategy reached a high level, but only the margin criterion of level 7 is
met. The other three need evidence that the run did not produce: a rebalancing test, a capacity estimate and a sealed
holdout.
\item \emph{ARL 6, depth 4/4} (hypothetical). A deflated Sharpe ratio of at least 0.99, low probability of overfitting,
an audited trial count and a passed holdout. It is still below any strategy at level 7.
\end{itemize}

\subsection{Principles of honest use}

The rules file states seven principles:
\begin{enumerate}
\item thresholds are versioned, a change creates a new version, and every certificate records the version used;
\item the sealed holdout is opened once;
\item trials are counted for the whole project, not for a single idea;
\item paper, testnet and simulated fills are not live and cap the level at 8;
\item missing evidence is shown and never counted;
\item depth never raises the level;
\item the scale is not a recommendation to trade and says nothing about future returns.
\end{enumerate}

\section{Worked example}\label{sec:example}

Table~\ref{tab:example} shows five strategies on artificial daily data in which the truth is known. The short program
\texttt{arl\_examples.py}, published with the sources of this paper, produces these results with Monte-Neo 0.57.0. All
runs charge the default costs of 10~bps per side. These are single illustrative cases; the rates over many cases are in
Sections~\ref{sec:results} and~\ref{sec:confirm}.

\begin{table}[ht]
\centering\small
\begin{tabularx}{\linewidth}{@{}Lrrrll@{}}
\toprule
Case & Sharpe & DSR & Trials & Result & Why it stops \\
\midrule
A. Uses tomorrow's price, 5 years & 22.13 & 1.000 & 1 & ARL 1, depth 2/4 & look-ahead \\
B. Best of 100 random signals, search not recorded, 3 years & 0.92 & 0.940 & 1 & ARL 3, depth 2/4 & no trial evidence \\
C. The same data, 100 trials declared, with a ledger & 0.92 & 0.164 & 100 & ARL 3, depth 2/4 & DSR below 0.8 \\
D. Real edge, 1 year & 1.90 & 0.461 & 20 & ARL 3, depth 2/4 & DSR below 0.8 \\
E. Real edge, 10 years & 1.98 & 1.000 & 20 & ARL 7, depth 1/4 & level 8 not run \\
\bottomrule
\end{tabularx}
\caption{Executed examples (net annual Sharpe ratio; DSR is the deflated Sharpe ratio).}\label{tab:example}
\end{table}

\begin{itemize}
\item \textbf{A} shows why the order of levels matters. A Sharpe ratio of 22 is worthless here: the look-ahead probes
fail, so the strategy stays at level 1.
\item \textbf{B} and \textbf{C} are the same random signal. Reported as one attempt, its deflated Sharpe ratio of 0.940
would pass the level-4 threshold; only the missing trial record stops it. Declaring the 100 attempts lowers the value to
0.164.
\item \textbf{D} and \textbf{E} share the same kind of true edge. One year is not enough to separate it from luck after
20 trials; ten years are. The level of D is a statement about the sample, not a proof that there is no edge.
\item \textbf{E} stops at level 7 because no book was supplied. ``Not assessed'' is different from ``failed''.
\end{itemize}

\section{Calibration design}\label{sec:calib}

Calibration asks two questions of the scale: how often does it promote what it should not, and how often does it
recognise what is real? Both can be answered only with cases whose truth is known.

\subsection{Three populations}

\paragraph{Population A: traps.} Each trap has a flaw and a maximum level that I set by its design. Single-asset traps:
\begin{itemize}
\item look-ahead by a shifted window, by a centred window, or by a z-score over the full sample (maximum 1);
\item costs larger than the gross edge (maximum 2);
\item phantom profits in pairs trading created by the way returns are booked \citep{shcherbyna2026phantom} (maximum 2);
\item hidden trials: the best of 100 or 1,000 random strategies reported as one trial (maximum 3);
\item a survivor universe: only assets that survived and rose (maximum 4);
\item regime decay: an edge that vanishes halfway through the sample (maximum 6).
\end{itemize}
Thirteen further traps rank many assets at once (the \emph{alpha zoo}); I gave each of them the same maximum, 4:
\begin{enumerate}
\item momentum presented as a new signal;
\item steady gains that are lost when the market falls;
\item an edge that lives in three assets only;
\item a score that pays only on the first trading day of each month;
\item a short side in thinly traded small stocks;
\item six years of today's members with no asset added or removed;
\item assets held before they joined the universe;
\item Tokyo stocks scored with New York returns of the same date, which are not yet known;
\item an edge that disappears after its publication date;
\item a score unrelated to returns;
\item a near copy of a strategy already in the book;
\item an edge that disappears under slight noise in the inputs;
\item a claimed information ratio far above what the evidence allows.
\end{enumerate}

\paragraph{Population B: no edge.} The best of $K$ random strategies or scores, with the true $K$ declared, on simulated
random walks and on real daily bars of 100 cryptocurrency pairs, 33 currency pairs, 12 stock indices and 7 commodity
series (maximum 3). A noise result can still be profitable by chance after costs, which is why level 3 is allowed and
level 4 is not.

\paragraph{Population C: real edges.} Single-asset strategies with a target net annual Sharpe ratio of 0.5, 1, 1.5 or 2
over 1, 3, 5 or 10 years; cross-sectional signals with four edge sizes (E1 smallest to E4); and honest ``twins'' of the
zoo traps, which have the same construction without the flaw. No maximum applies. Half of the single-asset cases declare
20 trials and half declare one.

\subsection{Tracks, measures and reproducibility}

\begin{itemize}
\item \textbf{Tracks.} Track V runs the verifier for single strategies, on daily bars except for the pairs-trading trap
(hourly returns only). Track X runs the report for rankings of many assets. Each case is scored on its own track.
Before the repair described in Section~\ref{sec:failures}, track V could not reach level 5; afterwards it reaches level
7. Level 8 needs a book, which only track X supplies.
\item \textbf{Measures.} For A and B I report the false promotion rate; for C, the power at each level. Intervals are
Wilson 95\,\% intervals and treat cases as independent.
\item \textbf{Reproducibility.} Every case is determined by its index and a base seed (real-data cases also by the file
of bars), so a larger run extends a smaller one.
\item \textbf{A limitation of development.} In the three development runs (base seed 20261119) the zoo cases draw their
data from the seed modulo 1,000. Of their 57,243 cases only 22,781 are distinct, so the development intervals for track
X in populations A and C are too narrow, and the three runs of Table~\ref{tab:calib} share every seed. The confirmatory
run uses new seeds and draws zoo data from the full case seed (the seed of the placebo permutations remains the case
seed modulo 1,000).
\end{itemize}

A confirmatory or development run holds 350,000 cases: 84,000 traps (24\,\%), 77,000 no-edge cases (22\,\%) and 189,000
cases of population C (54\,\%).

\subsection{Declared evidence}\label{sec:declared}

Part of the evidence is declared by the user, so the calibration must decide what each kind of case declares. In the
development run of the final rules:
\begin{itemize}
\item every case except the hidden-trials trap and the pairs-trading trap keeps a new trial ledger, as an honest
researcher would; the hidden-trials trap keeps none, as a researcher who hides trials would not;
\item every cross-sectional case except the zoo traps declares a fixed membership file, including the no-edge scores;
\item twelve of the thirteen zoo traps declare no membership file; the membership-window trap carries its own.
\end{itemize}
These choices decide part of the result. The hidden-trials trap is capped by construction, and 2,322 zoo-trap cases
(6.6\,\%) are stopped at level 5 only by the missing membership file; 1,196 of them are the survivor-universe trap, whose
flaw this check is meant to catch. The confirmatory run therefore measures both this design and one in which traps
receive the same declarations as honest cases.

\section{Development results}\label{sec:results}

\subsection{Two structural failures}\label{sec:failures}

\paragraph{Failure 1: hidden trials.} A full run of an early version of the second rules left 10.9\,\% of traps falsely
promoted. Ninety per cent of these came from one trap: the best of 100 or 1,000 random strategies reported as a single
trial reached level 4 in 92.7\,\% of its cases (the rest came from the cost-omission trap,
Section~\ref{sec:third}). No threshold can solve this. The selected series is, by construction, indistinguishable from
a series that one honest attempt produced with a lucky draw; no test on that series alone can see the discarded
variants.

The remedy is to ask for evidence about the search itself: levels 4 and above are capped at 3 unless the number of
trials is evidenced. This moves the burden from detection to documentation, and its limits must be stated plainly. In
the calibration the hidden-trials trap was run without a ledger and every other case, except the pairs-trading trap,
with a new one, so its result (0 of 8,912 at level 4 or above) holds by construction. The check verifies the chain of
hashes in the ledger, not when the ledger was opened or how many entries it should hold. When 40 hidden-trials cases
were rerun with a new, empty ledger, 18 of them (45\,\%, Wilson 95\,\% interval 31--60\,\%) reached level 4 again. The
ceiling stops a researcher who keeps no record, not one who opens a new ledger after the search.

\paragraph{Failure 2: no path for single strategies.} Under the first rules and the early second rules, no planted
single-asset edge reached level 5 (0 of 112,000). The two checks of level 5, residual alpha and hidden short volatility,
were computed only for rankings of many assets and for carry strategies. The verifier for single strategies now computes
both against the asset the strategy trades. These two rows warn at most, so a strategy whose return is only market
exposure stops at level 4 with warnings rather than being reported as a broken backtest.

\subsection{A third change and its cost}\label{sec:third}

The final rules contain a third change that is not structural: on the single-strategy track, level 3 also requires the
break-even cost to be at least 3 times the modelled cost per side (the earlier floor was 2). It removed the false
promotion of the cost-omission trap (900 of 8,912 cases, 10.1\,\%, to none), but it also lowered the share of planted
single-asset edges reaching level 3 from 90.7\,\% to 57.9\,\%. For a planted net Sharpe ratio of 0.5 the share reaching
level 3 now \emph{falls} as the sample grows (29.9, 19.3, 12.7 and 5.4\,\% for 1, 3, 5 and 10 years), because a longer
sample measures the thin margin more precisely. Level 3 under these rules therefore certifies a margin over costs, not
only a profit after costs.

The second repair had a cost too. With a path above level 4 open, the regime-decay trap, whose edge vanishes in the
second half of the sample, reached level 7, above its maximum of 6, in 7.0\,\% of cases (620 of 8,904; 17.6\,\% with ten
years of data). Every remaining false promotion of a trap in development comes from this trap.

\subsection{Three versions of the rules}

\begin{table}[ht]
\centering\small
\begin{tabular}{@{}lccc@{}}
\toprule
 & Rules 1 & Rules 2a & Rules 2c\\
\midrule
Cases & 350,000 & 350,000 & 350,000 \\
False promotion, A (traps) & 27.6\,\% & 10.9\,\% & 0.74\,\% [0.68, 0.80] \\
False promotion, B (no edge) & 33.8\,\% & 1.0\,\% & 0.77\,\% [0.71, 0.84] \\
\quad of which real market bars & 43.8\,\% & 1.9\,\% & 1.73\,\% [1.57, 1.90] \\
Hidden-trials trap at level 4 or above & 96.9\,\% & 92.7\,\% & 0 of 8,912$^{a}$ \\
Regime-decay trap above its maximum (6) & 0\,\% & 0\,\% & 7.0\,\% [6.5, 7.5] \\
Planted single-asset edge, level 3 or above & 90.7\,\% & 90.7\,\% & 57.9\,\% [57.6, 58.2] \\
Planted single-asset edge, level 6 or above & 0\,\% & 0\,\% & 28.9\,\% [28.7, 29.2] \\
Planted cross-sectional edge, level 6 or above & 54.9\,\% & 52.8\,\% & 52.8\,\% [52.4, 53.2] \\
\bottomrule
\end{tabular}
\caption{Three versions of the rules on the same 350,000 cases. Rules 1: \texttt{arl-rules/1}. Rules 2a: the deflated
Sharpe floor at level 4 and the blockers at level 5 (commit \texttt{766213e}). Rules 2c: Rules 2a plus the
trial-evidence ceiling, attribution for single strategies and the cost margin of 3 at level 3 (commit \texttt{91171e3};
unchanged in release \texttt{v0.57.0}, commit \texttt{cd88c80}). An intermediate version 2b is not shown. Brackets:
Wilson 95\,\% interval. $^{a}$By construction (Section~\ref{sec:failures}). Cross-sectional edges: E1--E4 only, without
the honest twins.}\label{tab:calib}
\end{table}

The step from the first to the second rules removed most false promotion of noise. The later changes removed the false
promotion of the hidden-trials and cost-omission traps and gave single strategies a path above level 4, without changing
the power on cross-sectional signals, at the price of a new false promotion of the regime-decay trap and lower
single-asset power at level 3. All three versions were evaluated on the cases on which their thresholds were chosen, so
these rates are in-sample estimates. Section~\ref{sec:confirm} tests the last column on new cases.

\paragraph{Versions and names.} Rules 2a and 2c both carried the name \texttt{arl-rules/2} while the rules were still
under development. Rules 2a existed only as a development snapshot in the source repository; no packaged release
contains it. The first release with \texttt{arl-rules/2} is \texttt{v0.57.0}, and it contains Rules 2c. The snapshots
are identified here by their commits.

\paragraph{Real market bars} are the harder test with no edge. Under Rules 2c the best random strategy went above level 3
in:
\begin{itemize}
\item cryptocurrency pairs: 3.8\,\% (389 of 10,143);
\item commodities: 2.0\,\% (22 of 1,115);
\item stock indices: 0.4\,\% (8 of 1,889);
\item currency pairs: 0.04\,\% (5 of 11,349).
\end{itemize}
The classes are not directly comparable: all ten-year windows are currency pairs, and every window is annualised with
365 bars a year although markets that close at weekends have fewer.

\section{Confirmatory run}\label{sec:confirm}

\subsection{Design}

Before the run, I deposited a plan with four predictions, the case generator, the analysis script and their hashes at
Zenodo \citep{shcherbyna2026prereg}. The record is time-stamped 27 November 2026, 20:25 UTC; the launcher recorded the
first start at 21:10 UTC. The verifier was frozen at release \texttt{v0.57.0} (commit \texttt{cd88c80}), and the base
seed, 20261128, had not been run before. The 350,000 cases follow the development schedule and were run in two
\emph{arms} on the same seeds:
\begin{itemize}
\item \textbf{Silent arm} (the development design): the hidden-trials trap keeps no ledger, and the cross-sectional traps
are given no added membership file.
\item \textbf{Parity arm} (primary): every case receives a new trial ledger, and every cross-sectional case that declares
no membership file receives one listing a fixed universe. The pairs-trading trap, which has returns only, takes neither.
\end{itemize}
A new ledger holds no earlier trial; for the hidden-trials trap it declares one trial although 100 or 1,000 variants were
searched. For a trap, the parity arm models a user who supplies the cheapest declarations that the rules accept. Only the
43,912 cases whose inputs differ between the arms were run again for the silent arm.

\paragraph{Pilot.} Before writing the plan I had run 1,500 cases with another base seed (20261127) in both arms to test
the software, and I had seen their parity-arm results: 15 of 40 hidden-trials cases reached level 4, and some
cross-sectional traps were promoted. Predictions P2 and P3 are therefore informed by a pilot and are not blind.

\paragraph{Predictions} (bounds apply to point estimates):
\begin{description}
\term{P1.} In the silent arm, eight rates lie within 1.0 percentage point of their development values (1.5 points for
the regime-decay trap), and the hidden-trials trap never reaches level 4.
\term{P2.} In the parity arm, the hidden-trials trap reaches level 4 or above in 25--55\,\% of cases.
\term{P3.} In the parity arm, the pooled false promotion of traps is 4--14\,\%.
\term{P4.} A separate full run in the silent configuration of the first 3,500 cases gives the same level, depth,
blocking level, blocking rule and ceilings as the arms it should match.
\end{description}

\subsection{Results}

All four predictions were met:
\begin{description}
\term{P1} met: the largest difference was 0.56 points (regime-decay trap); all others were at most 0.26 points.
\term{P2} met: 46.9\,\%.
\term{P3} met: 8.29\,\%.
\term{P4} met: 3,062 cases were compared with the parity arm and 438 with the silent arm, with no difference.
\end{description}
All 350,000 cases completed without a software error, and none had to be retried. One departure concerns logistics
only: after a network failure the verifier was cloned from a local copy of the public repository; before starting, the
launcher checked the commit, the unchanged checkout, and the hashes of the deposited files and of the file of bars.

\begin{table}[ht]
\centering\small
\begin{tabular}{@{}lccc@{}}
\toprule
 & Development & Confirmatory, silent & Confirmatory, parity\\
\midrule
False promotion, A (traps) & 0.74\,\% & 0.80\,\% [0.74, 0.86] & 8.29\,\% [8.10, 8.48] \\
False promotion, B (no edge) & 0.77\,\% & 0.82\,\% [0.76, 0.89] & 0.82\,\% [0.76, 0.89] \\
\quad of which real market bars & 1.73\,\% & 1.81\,\% [1.65, 1.98] & 1.81\,\% [1.65, 1.98] \\
Hidden-trials trap at level 4 or above & 0\,\% & 0\,\% & 46.9\,\% [45.8, 47.9] \\
Regime-decay trap above its maximum & 7.0\,\% & 7.5\,\% [7.0, 8.1] & 7.5\,\% [7.0, 8.1] \\
Planted single-asset edge, level 3 or above & 57.9\,\% & 58.2\,\% [57.9, 58.5] & 58.2\,\% [57.9, 58.5] \\
Planted single-asset edge, level 6 or above & 28.9\,\% & 29.0\,\% [28.8, 29.3] & 29.0\,\% [28.8, 29.3] \\
Planted cross-sectional edge, level 3 or above & 79.6\,\% & 79.7\,\% [79.3, 80.0] & 79.7\,\% [79.3, 80.0] \\
Planted cross-sectional edge, level 6 or above & 52.8\,\% & 52.8\,\% [52.4, 53.2] & 52.8\,\% [52.4, 53.2] \\
\bottomrule
\end{tabular}
\caption{Frozen final rules on 350,000 development cases and 350,000 new cases. Brackets: Wilson 95\,\% interval.
No-edge and planted-edge cases have the same inputs in both arms. The zeros for the hidden-trials trap in the first two
columns hold by construction.}\label{tab:confirm}
\end{table}

\subsection{Where traps pass when they declare evidence}\label{sec:parity}

Table~\ref{tab:traps} breaks down the 6,962 falsely promoted traps of the parity arm. The explanations come from the
check results of the individual cases; this analysis was not pre-registered.

\begin{table}[ht]
\centering\small
\begin{tabularx}{\linewidth}{@{}p{0.2\linewidth}rrrL@{}}
\toprule
Trap & Max. & Above maximum & Highest & What the check results show \\
\midrule
Hidden trials & 3 & 4,177 / 8,912 & 7 & the ledger check cannot see trials that the ledger does not record \\
Survivor universe (zoo) & 4 & 1,169 / 2,692 & 7 & the false universe declaration is accepted; the survivorship ceiling exists only on the single-strategy track \\
Regime decay & 6 & 670 / 8,904 & 7 & period consistency passes; holdout and regime warnings do not affect the level \\
Inflated claim & 4 & 606 / 2,692 & 7 & the breadth check warns but does not affect the level \\
Copy of the book & 4 & 318 / 2,692 & 8 & the residual-against-book check passes; the drawdown-overlap warning does not affect the level \\
Decay after discovery & 4 & 21 / 2,692 & 7 & the decay check does not warn \\
Cost omission & 2 & 1 / 8,912 & 6 & a single case \\
Fourteen other traps & 1--4 & 0 / 46,504 & -- & -- \\
\bottomrule
\end{tabularx}
\caption{False promotion by trap in the parity arm. ``Highest'' is the highest level among the falsely promoted
cases.}\label{tab:traps}
\end{table}

Three kinds of weakness appear.
\begin{enumerate}
\item \textbf{The scale trusts declarations.} The hidden-trials and the survivor-universe traps account for 5,346 of the
6,962 cases. The hidden-trials trap reached level 4 in 3,682 cases, level 6 in 6 and level 7 in 489. The survivor case
also repeats the second structural failure: the ceiling at level 4 for a universe of survivors is computed only on the
single-strategy track, where all 4,452 survivor cases carried it and none went above level 3. On the cross-sectional
track, declaring a membership file removes the only check aimed at this flaw.
\item \textbf{Some warnings are not connected to the levels.}
\begin{itemize}
\item The check of the claimed against the attainable information ratio (\chk{breadth\_consistency}) warned in all
2,692 inflated-claim cases, while the blocker meant for claims (\chk{claim\_consistency}) is not computed by the
cross-sectional report.
\item The check of overlap in bad periods with the book (\chk{drawdown\_overlap}) warned in all 318 promoted copies of
the book.
\item Among the 670 promoted regime-decay cases, \chk{holdout\_consistency} warned in 259 and \chk{regime\_stability}
in 71 (at least one of the two in 286).
\end{itemize}
I have not measured what connecting these checks would cost on honest cases.
\item \textbf{Some detectors miss.} The residual-against-book check passed in all 318 promoted copies (the marginal
Sharpe check passed in 316). The decay check, which warned in 2,638 of 2,692 cases, did not warn in any of the 21
promoted ones.
\end{enumerate}
The silent arm hides all of these except the regime-decay promotions: its hidden-trials trap is capped for lack of a
ledger, and 2,285 of the 35,000 zoo-trap cases are stopped at level 5 by the missing membership file alone; 2,114 of
them exceed their maximum once the file is supplied. I did not change the rules after the run.

On real market bars the best random strategy went above level 3 in the new cases in:
\begin{itemize}
\item cryptocurrency pairs: 4.0\,\% (413 of 10,220);
\item commodities: 1.4\,\% (15 of 1,070);
\item stock indices: 0.32\,\% (6 of 1,868);
\item currency pairs: 0.08\,\% (9 of 11,338).
\end{itemize}
The series are those of development; only the windows and the random signals are new.

\subsection{Power: how much data a real edge needs}

\begin{table}[ht]
\centering\small
\begin{tabular}{@{}lcccc@{}}
\toprule
Single-asset edge (annual Sharpe) & 1 year & 3 years & 5 years & 10 years\\
\midrule
0.5 & 0.3\,\% & 1.4\,\% & 2.0\,\% & 4.2\,\% \\
1.0 & 1.5\,\% & 5.8\,\% & 14.0\,\% & 38.0\,\% \\
1.5 & 4.1\,\% & 20.6\,\% & 48.3\,\% & 87.0\,\% \\
2.0 & 8.0\,\% & 47.6\,\% & 82.1\,\% & 99.6\,\% \\
\midrule
Cross-sectional edge & E1 & E2 & E3 & E4 \\
\midrule
Level 3 or above & 27.8\,\% & 91.6\,\% & 99.4\,\% & $>$99.9\,\% \\
Level 6 or above & 0.8\,\% & 39.8\,\% & 79.7\,\% & 90.8\,\% \\
\bottomrule
\end{tabular}
\caption{Power in the confirmatory run: share of planted edges reaching level 6 (single-asset track, 7,000 cases per
cell; Wilson 95\,\% intervals within $\pm$1.2 points) and levels 3 and 6 (cross-sectional track, about 13,700 cases per
edge size).}\label{tab:power}
\end{table}

Table~\ref{tab:power} shows that power at level 6 rises with the size of the edge and the length of the sample:
\begin{itemize}
\item a true annual Sharpe ratio of 2 reaches level 6 in 99.6\,\% of ten-year samples but in only 8.0\,\% of one-year
samples;
\item a true Sharpe ratio of 0.5 almost never reaches it (at most 4.2\,\%);
\item on the cross-sectional track the smallest edge reached level 6 in 0.8\,\% of cases and the largest in about
91\,\%;
\item at level 3 the weakest single-asset edge behaves in the opposite way, falling from 30.4\,\% with one year to
5.9\,\% with ten, because of the cost margin of Section~\ref{sec:third}.
\end{itemize}
The practical message for a researcher with a short backtest is that a level below 6 often says something about the
sample rather than about the edge. A weak edge is not promoted early; it needs more data, and the scale says which
evidence is missing.

\section{Using and reporting ARL}\label{sec:use}

I propose ARL as an open way of reporting the evidence behind a backtest. It becomes a standard only if others adopt it,
test it and improve it. For results to be comparable, a report of an ARL level should state:
\begin{enumerate}
\item the level and depth, for example ``ARL 6, depth 3/4'';
\item the rules version and the software release;
\item the kind of strategy and the track used;
\item the trial evidence: its type, the number of trials, and when the record was started;
\item every declaration that was not verified, such as a membership file or the source of live results;
\item the checks that warned, and the levels that were not assessed;
\item the hash of the data, so that the certificate can be recomputed.
\end{enumerate}
An example of a complete one-line statement is: ``ARL 4, depth 2/4 (\texttt{arl-rules/2}, Monte-Neo 0.57.0); 35 trials in
a ledger opened before the first test; level 5 not assessed (no attribution run).''

\paragraph{Advice to a researcher.} Open the trial ledger before the first test, not after the best one. Record every
variant. Fix a holdout and open it once. Treat the first level that is not met as the next task.

\paragraph{Advice to a reader or investor.} Ask when the trial record was started and what it contains. Ask who declared
the universe. Do not read a level as a promise of profit. Do not read a testnet or demo record as live evidence.

\paragraph{Advice to builders of verification tools.} Report the false promotion rate both for traps that declare nothing
and for traps that declare everything an honest case declares. The difference shows how much of the scale rests on
trust.

\section{Implementation and reproducibility}\label{sec:impl}

\begin{itemize}
\item The rules are one JSON file (\texttt{arl\_rules\_v2.json}) read by the verifier. Documentation tables are generated
from it, and a test fails when they differ.
\item Earlier rule versions remain loadable, so a certificate can be recomputed with the version it records.
\item Each certificate stores the rules version, the settings and a hash of the data; a recheck with the same software
version reproduces it.
\item The trial ledger is an append-only, hash-chained file; an edited ledger is reported and is not accepted as
evidence.
\item The generator deposited with this paper differs from the one used in development only in its configuration; with
the development settings it reproduced the level, depth, blocking rule, ceilings and verdict of the first 3,000
development cases of the final rules.
\end{itemize}
Monte-Neo is written by the author of this paper, who therefore designed both the checks and the traps. No independent
party has yet tested the scale.

Monte-Neo also labels each verified strategy with one of 24 main categories (how the signal is formed, for example
time-series trend, cross-sectional momentum, mean reversion, statistical arbitrage or carry), an asset class and a
structure, derived by twelve ordered rules (\texttt{catalog-rules/1}), so that levels can be compared among strategies of
the same kind.

\section{Limitations}

\begin{itemize}
\item \textbf{Mostly synthetic data.} Planted edges are stationary, whereas real edges decay \citep{mclean2016}.
\item \textbf{Author-chosen traps.} The traps are the ones I thought of; a new kind of error may pass. The confirmatory
run uses new random cases from the same generator: it removes the selection of thresholds on the measured cases, but not
the dependence on my choice of case types. Its real-data cases reuse the development series.
\item \textbf{Provisional thresholds.} They remain provisional until a collection of real strategies with known outcomes
is available.
\item \textbf{Unverified declarations.} Any intact trial ledger meets the trial-evidence requirement, including one
opened after the search. It becomes evidence only when it is fixed before the search, for example by a time-stamped
pre-registration or a hash deposited with a third party, which the rules do not require. A membership file is accepted
as written.
\item \textbf{The parity rate is not a field rate.} The 8.29\,\% describes the cheapest accepted declarations for my mix
of traps, not what to expect from real users. The hidden-trials trap alone supplies 4,177 of the 6,962 cases; without it
the rate is 3.71\,\% (2,785 of 75,088).
\item \textbf{Depth is not calibrated.} The out-of-sample criterion was not assessed in any calibration case. One
criterion is also hard to meet by construction: with 200 permutations or shifts the smallest possible p-value is about
0.005, so D1 of level 4 cannot be met with three or more trials.
\item \textbf{Coverage.} Intervals treat cases as independent. Level 8 can be reached only on the cross-sectional track,
and no case reached level 9, so the live requirements are untested here.
\item \textbf{Returns alone are not enough.} They cannot test for look-ahead or costs, and the scale leaves those levels
not assessed.
\item \textbf{Evidence is not profit.} ARL measures how well a claim is supported, not what will be earned.
\end{itemize}

\section{Conclusions}

\paragraph{What this paper is for.} It offers a common language for one question that every user of a backtest should
ask: how much has actually been shown? A level between 1 and 9, with a depth, answers it in a form that can be checked.

\paragraph{What was found.}
\begin{enumerate}
\item The first serious failures of the scale were failures of structure, not of thresholds. Hidden trials cannot be
detected from a result and must be documented; a check that exists on one track but not on another silently blocks real
strategies.
\item On new cases and with rules frozen in advance, no-edge strategies exceeded their maximum in 0.82\,\% of cases, and
the development rates were reproduced.
\item The ability to recognise a real edge depends strongly on how much data exist.
\item A low rate of fooled traps can be an artefact of the test design. When traps declared what honest cases declare,
the rate rose from 0.80\,\% to 8.29\,\%.
\item A level is therefore conditional: it is as trustworthy as the records and declarations behind it.
\end{enumerate}

\paragraph{What follows.} I recommend that any readiness scale for backtests report both rates. For ARL the next steps
follow from Table~\ref{tab:traps}: fix the trial ledger and the universe declaration before the search; compute the
survivorship ceiling on both tracks; connect the warning checks to the levels where their cost on honest cases allows;
strengthen the book and decay detectors; and test the next version on new cases and on traps built by others.

\section*{Availability and declarations}

\begin{itemize}
\item \textbf{Software.} Monte-Neo and the rules file are open source (MIT licence) at
\url{https://github.com/NeoZorK/Monte-Neo}; this study uses release \texttt{v0.57.0} and rules \texttt{arl-rules/2}. The
sources of this paper and the example program are in the folder \texttt{paper/arl-standard} of the same repository.
\item \textbf{Pre-registration.} The plan, case generator, analysis script and launcher are deposited at Zenodo
(\href{https://doi.org/10.5281/zenodo.24603671}{doi:10.5281/zenodo.24603671}).
\item \textbf{Results.} The results of every case of the development and confirmatory runs (level, depth, blocking rule,
ceilings, status of every check and principal metrics) are deposited as a later version of the same record (all
versions: \href{https://doi.org/10.5281/zenodo.24603670}{doi:10.5281/zenodo.24603670}). Further numerical details of the
checks are omitted for size; the deposited generator recreates them from the case seeds.
\item \textbf{Real data.} The daily bars used for part of population B are listed there by symbol, asset class, date
range and hash. They are not redistributed.
\item \textbf{Authorship of the scale.} The Alpha Readiness Level was devised by the author.
\item \textbf{Funding and interests.} The author received no funding. The author develops the software evaluated in this
study.
\item Nothing in this paper is investment advice.
\end{itemize}

\begin{thebibliography}{99}
\bibitem[Arnott et al.(2019)]{arnott2019backtesting} Arnott, R., Harvey, C.~R., and Markowitz, H. (2019). A backtesting
protocol in the era of machine learning. \emph{Journal of Financial Data Science}, 1(1), 64--74. \href{https://doi.org/10.3905/jfds.2019.1.064}{doi:10.3905/jfds.2019.1.064}.
\bibitem[Bailey and L\'opez de Prado(2014)]{bailey2014deflated} Bailey, D.~H., and L\'opez de Prado, M. (2014). The deflated
Sharpe ratio: correcting for selection bias, backtest overfitting, and non-normality. \emph{Journal of Portfolio
Management}, 40(5), 94--107. \href{https://doi.org/10.3905/jpm.2014.40.5.094}{doi:10.3905/jpm.2014.40.5.094}.
\bibitem[Bailey et al.(2017)]{bailey2017pbo} Bailey, D.~H., Borwein, J., L\'opez de Prado, M., and Zhu, Q.~J. (2017). The
probability of backtest overfitting. \emph{Journal of Computational Finance}, 20(4), 39--69. \href{https://doi.org/10.21314/JCF.2016.322}{doi:10.21314/JCF.2016.322}.
\bibitem[Bailey et al.(2014)]{bailey2014pseudo} Bailey, D.~H., Borwein, J.~M., L\'opez de Prado, M., and Zhu, Q.~J. (2014).
Pseudo-mathematics and financial charlatanism: the effects of backtest overfitting on out-of-sample performance.
\emph{Notices of the American Mathematical Society}, 61(5), 458--471. \href{https://doi.org/10.1090/noti1105}{doi:10.1090/noti1105}.
\bibitem[Fabozzi and L\'opez de Prado(2018)]{fabozzi2018honest} Fabozzi, F.~J., and L\'opez de Prado, M. (2018). Being
honest in backtest reporting: a template for disclosing multiple tests. \emph{Journal of Portfolio Management}, 45(1),
141--147. \href{https://doi.org/10.3905/jpm.2018.45.1.141}{doi:10.3905/jpm.2018.45.1.141}.
\bibitem[Getmansky et al.(2004)]{getmansky2004} Getmansky, M., Lo, A.~W., and Makarov, I. (2004). An econometric model of
serial correlation and illiquidity in hedge fund returns. \emph{Journal of Financial Economics}, 74(3), 529--609. \href{https://doi.org/10.1016/j.jfineco.2004.04.001}{doi:10.1016/j.jfineco.2004.04.001}.
\bibitem[Hansen(2005)]{hansen2005spa} Hansen, P.~R. (2005). A test for superior predictive ability. \emph{Journal of
Business and Economic Statistics}, 23(4), 365--380. \href{https://doi.org/10.1198/073500105000000063}{doi:10.1198/073500105000000063}.
\bibitem[Harvey et al.(2016)]{harvey2016cross} Harvey, C.~R., Liu, Y., and Zhu, H. (2016). \ldots and the cross-section of
expected returns. \emph{Review of Financial Studies}, 29(1), 5--68. \href{https://doi.org/10.1093/rfs/hhv059}{doi:10.1093/rfs/hhv059}.
\bibitem[Lavin et al.(2022)]{lavin2022mltrl} Lavin, A., Gilligan-Lee, C.~M., Visnjic, A., et al. (2022). Technology
readiness levels for machine learning systems. \emph{Nature Communications}, 13, 6039. \href{https://doi.org/10.1038/s41467-022-33128-9}{doi:10.1038/s41467-022-33128-9}.
\bibitem[Lo(2002)]{lo2002sharpe} Lo, A.~W. (2002). The statistics of Sharpe ratios. \emph{Financial Analysts Journal},
58(4), 36--52. \href{https://doi.org/10.2469/faj.v58.n4.2453}{doi:10.2469/faj.v58.n4.2453}.
\bibitem[Mankins(1995)]{mankins1995trl} Mankins, J.~C. (1995). \emph{Technology Readiness Levels: A White Paper}. NASA,
Office of Space Access and Technology.
\bibitem[McLean and Pontiff(2016)]{mclean2016} McLean, R.~D., and Pontiff, J. (2016). Does academic research destroy stock
return predictability? \emph{Journal of Finance}, 71(1), 5--32. \href{https://doi.org/10.1111/jofi.12365}{doi:10.1111/jofi.12365}.
\bibitem[Nosek et al.(2018)]{nosek2018prereg} Nosek, B.~A., Ebersole, C.~R., DeHaven, A.~C., and Mellor, D.~T. (2018). The
preregistration revolution. \emph{Proceedings of the National Academy of Sciences}, 115(11), 2600--2606. \href{https://doi.org/10.1073/pnas.1708274114}{doi:10.1073/pnas.1708274114}.
\bibitem[Shcherbyna(2026a)]{shcherbyna2026phantom} Shcherbyna, R. (2026a). Phantom mean reversion: how rolling hedge ratios
manufacture profits in pairs-trading backtests. Preprint, Zenodo, \href{https://doi.org/10.5281/zenodo.24561375}{doi:10.5281/zenodo.24561375}.
\bibitem[Shcherbyna(2026b)]{shcherbyna2026prereg} Shcherbyna, R. (2026b). Pre-registration: confirmatory calibration of
the Alpha Readiness Level. Zenodo, \href{https://doi.org/10.5281/zenodo.24603671}{doi:10.5281/zenodo.24603671}.
\bibitem[White(2000)]{white2000reality} White, H. (2000). A reality check for data snooping. \emph{Econometrica}, 68(5),
1097--1126. \href{https://doi.org/10.1111/1468-0262.00152}{doi:10.1111/1468-0262.00152}.
\bibitem[Wilson(1927)]{wilson1927} Wilson, E.~B. (1927). Probable inference, the law of succession, and statistical
inference. \emph{Journal of the American Statistical Association}, 22(158), 209--212. \href{https://doi.org/10.1080/01621459.1927.10502953}{doi:10.1080/01621459.1927.10502953}.
\end{thebibliography}

\end{document}
